%% LaTeX-Beamer template for KIT design
%% by Erik Burger, Christian Hammer
%% title picture by Klaus Krogmann
%%
%% version 2.1
%%
%% mostly compatible to KIT corporate design v2.0
%% http://intranet.kit.edu/gestaltungsrichtlinien.php
%%
%% Problems, bugs and comments to
%% burger@kit.edu

\documentclass[18pt]{beamer}
\usepackage[utf8]{inputenc}
\usepackage{listings}
%% SLIDE FORMAT

% use 'beamerthemekit' for standard 4:3 ratio
% for widescreen slides (16:9), use 'beamerthemekitwide'

\usepackage{templates/beamerthemekit}
% \usepackage{templates/beamerthemekitwide}

%% TITLE PICTURE

% if a custom picture is to be used on the title page, copy it into the 'logos'
% directory, in the line below, replace 'mypicture' with the 
% filename (without extension) and uncomment the following line
% (picture proportions: 63 : 20 for standard, 169 : 40 for wide
% *.eps format if you use latex+dvips+ps2pdf, 
% *.jpg/*.png/*.pdf if you use pdflatex)

\titleimage{../media/banner2v2}

%% TITLE LOGO

% for a custom logo on the front page, copy your file into the 'logos'
% directory, insert the filename in the line below and uncomment it

\titlelogo{logo}

% (*.eps format if you use latex+dvips+ps2pdf,
% *.jpg/*.png/*.pdf if you use pdflatex)

%% TikZ INTEGRATION

% use these packages for PCM symbols and UML classes
% \usepackage{templates/tikzkit}
% \usepackage{templates/tikzuml}

% speech bubbles
\usepackage{tikz}
\usetikzlibrary{shapes.callouts}
\usetikzlibrary{arrows,positioning}
\usetikzlibrary{calc}

% rotations
\usepackage{rotating}

% the presentation starts here

\title[Croggle]{Croggle}
\subtitle{Spielend die Denkweise des Programmierens lernen}
\author{PSE 2013/2014}

\institute{Lukas Böhm $\bullet$ Tobias Hornberger $\bullet$ Jonas Mehlhaus $\bullet$ Iris Mehrbrodt $\bullet$ Vincent Schüßler $\bullet$ Lena Winter}

\beamertemplatenavigationsymbolsempty

\begin{document}

% change the following line to "ngerman" for German style date and logos
\selectlanguage{ngerman}

%title page
\begin{frame}
\titlepage
\end{frame}

%table of contents
\begin{frame}{Übersicht}
\tableofcontents
\end{frame}

%content
\section{Einführung}

\begin{frame}{Die Hauptdarsteller}
	\begin{center}
		\begin{tabular}{rl}
			\raisebox{-.5\height}{\includegraphics[height=2cm]{media/alligator_blank}}
			&
			\raisebox{-.5\height}{\includegraphics[height=2cm]{media/egg_blank}}
		\end{tabular}
	\end{center}
\end{frame}

\begin{frame}{Familien}
	\begin{center}
		\begin{tabular}{rccl}
			\multicolumn{4}{c}{
				\raisebox{-.5\height}{\includegraphics[height=2cm]{media/alligator_green}}
			}
			\\
			\multicolumn{2}{c}{
				\raisebox{-.5\height}{\includegraphics[height=1.5cm]{media/egg_green}}
			}
			&
			\multicolumn{2}{c}{
				\raisebox{-.5\height}{\includegraphics[height=1.5cm]{media/alligator_red}}
			}
			\\
			& &
			\raisebox{-.5\height}{\includegraphics[height=1cm]{media/egg_red}}
			&
			\raisebox{-.5\height}{\includegraphics[height=1cm]{media/egg_red}}
		\end{tabular}
	\end{center}
\end{frame}

\begin{frame}{Fressregeln: Beispiel 1}
	\begin{center}
		\begin{tabular}{rcccl}
			\multicolumn{2}{c}{
				\raisebox{-.5\height}{\includegraphics[height=1.25cm]{media/alligator_green}}
				\tikz[remember picture] \node (a) {\vphantom{X}};
				\begin{tikzpicture}[remember picture,overlay]
					\path<2> (a.east) ++(0,1) node[draw,anchor=west,ellipse callout,fill=white!50,opacity=1., callout absolute pointer={(a.mid)}]  {RAWR!};
				\end{tikzpicture}
			}
			& &
			\multicolumn{2}{c}{
				\raisebox{-.5\height}{\includegraphics[height=1.25cm]{media/alligator_yellow}}
				\tikz[remember picture] \node (b) {\vphantom{X}};
				\begin{tikzpicture}[remember picture,overlay]
					\path<2> (b.east) ++(0,1) node[draw,anchor=west,ellipse callout,fill=white!50,opacity=1., callout absolute pointer={(b.mid)}]  {HILFE!};
				\end{tikzpicture}
			}
			\\
			\multicolumn{2}{c}{
				\raisebox{-.5\height}{\includegraphics[height=1.25cm]{media/alligator_red}}
			}
			& &
			\multicolumn{2}{c}{
				\raisebox{-.5\height}{\includegraphics[height=0.75cm]{media/egg_yellow}}
			}
			\\
			\raisebox{-.5\height}{\includegraphics[height=.75cm]{media/egg_green}}
			&
			\raisebox{-.5\height}{\includegraphics[height=.75cm]{media/egg_red}}
		\end{tabular}
	\end{center}
\end{frame}

\begin{frame}{Fressregeln: Beispiel 1}
	\begin{center}
		\begin{picture}(320,80)
			\put(105,0){
				\tikz[remember picture] \node (a) {
					\begin{turn}{7}
						\includegraphics[height=1.25cm]{media/alligator_green}
					\end{turn}
				};
				\begin{tikzpicture}[remember picture,overlay]
					\path (a.north) ++(-1,1) node[draw,anchor=north,ellipse callout,fill=white!50,opacity=1., callout absolute pointer={(a.north)}]  {YUMMY!};
				\end{tikzpicture}
			}
			\put(180,20){
				\begin{rotate}{15}
					\parbox{2cm}{
						\tikz[remember picture] \node (b) {
							\includegraphics[height=1.25cm]{media/alligator_yellow}
						};
						\\
						\hspace*{1.3cm}
						\includegraphics[height=0.75cm]{media/egg_yellow}
					}
				\end{rotate}
				\begin{tikzpicture}[remember picture,overlay]
					\path (b.north) ++(-1,1) node[draw,anchor=north,ellipse callout,fill=white!50,opacity=1., callout absolute pointer={(b.north)}]  {AAAAAH!};
				\end{tikzpicture}
			}
		\end{picture}
		
		\begin{tabular}{lccc}
			\multicolumn{2}{l}{
				\raisebox{-.5\height}{\includegraphics[height=1.25cm]{media/alligator_red}}
			}
			\\
			\raisebox{-.5\height}{\includegraphics[height=.75cm]{media/egg_green}}
			&
			\raisebox{-.5\height}{\includegraphics[height=.75cm]{media/egg_red}}
		\end{tabular}
	\end{center}
\end{frame}

\begin{frame}{Fressregeln: Beispiel 1}
	\begin{center}
		\begin{tabular}{cc}
			\multicolumn{2}{c}{
				\tikz[remember picture] \node (a) {\vphantom{X}};
				\begin{turn}{180}
					\raisebox{-.5\height}{\includegraphics[height=1.25cm]{media/alligator_green}}
				\end{turn}
				\begin{tikzpicture}[remember picture,overlay]
					\path (a.north) ++(1,1) node[draw,anchor=east,ellipse callout,fill=white!50,opacity=1., callout absolute pointer={(a.mid)}]  {MIR IST SCHLECHT :(};
				\end{tikzpicture}
			}
			\\
			\multicolumn{2}{c}{
				\includegraphics[height=1.25cm]{media/alligator_red}
			}
			\\
			\includegraphics[height=.75cm]{media/egg_green}
			&
			\includegraphics[height=.75cm]{media/egg_red}
		\end{tabular}
	\end{center}
\end{frame}

\begin{frame}{Fressregeln: Beispiel 1}
	\begin{center}
		\begin{tabular}{cc}
			\multicolumn{2}{c}{
				\includegraphics[height=1.25cm]{media/alligator_red}
			}
			\\
			\tikz[remember picture] \node (a) {\includegraphics[height=.75cm]{media/egg_green}};
			\begin{tikzpicture}[remember picture,overlay]
				% add align option for multi line text
				\path (a.south) ++(0,-1) node[draw,anchor=east,ellipse callout,fill=white!50,opacity=1., callout absolute pointer={(a.south)},align=center]  {KRACK!\\[3mm] KRACK!};
			\end{tikzpicture}
			&
			\includegraphics[height=.75cm]{media/egg_red}
		\end{tabular}
	\end{center}
\end{frame}

\begin{frame}{Fressregeln: Beispiel 1}
	\begin{center}
		\begin{tabular}{ccc}
			\multicolumn{3}{c}{
				\includegraphics[height=1.25cm]{media/alligator_red}
			}
			\\
			\multicolumn{2}{c}{
				\includegraphics[height=1.25cm]{media/alligator_yellow}
				\tikz[remember picture] \node (a) {\vphantom{X}};
				\begin{tikzpicture}[remember picture,overlay]
					\path<2> (a.east) ++(0,-1) node[draw,anchor=west,ellipse callout,fill=white!50,opacity=1., callout absolute pointer={(a.mid)}]  {LECKER!};
				\end{tikzpicture}
			}
			&
			\includegraphics[height=.75cm]{media/egg_red}
			\\
			\multicolumn{2}{c}{
				\includegraphics[height=.75cm]{media/egg_yellow}
			}
		\end{tabular}
	\end{center}
\end{frame}

\begin{frame}{Fressregeln: Beispiel 1}
	\begin{center}
		\includegraphics[height=1.25cm]{media/alligator_red}
		\\
		\includegraphics[height=1.25cm]{media/alligator_yellow}
		\\
		\includegraphics[height=.75cm]{media/egg_yellow}
	\end{center}
\end{frame}

\begin{frame}{Fressregeln: Beispiel 1}
	\begin{center}
		\includegraphics[height=1.25cm]{media/alligator_red}
		\\
		\tikz[remember picture] \node (a) {
			\begin{turn}{180}
				\includegraphics[height=1.25cm]{media/alligator_yellow}
			\end{turn}
		};
		\begin{tikzpicture}[remember picture,overlay]
			\path (a.west) ++(0,-1) node[draw,anchor=east,ellipse callout,fill=white!50,opacity=1., callout absolute pointer={(a.west)}]  {AUTSCH! :(};
		\end{tikzpicture}
		\\
		\includegraphics[height=.75cm]{media/egg_yellow}
	\end{center}
\end{frame}

\begin{frame}{Fressregeln: Beispiel 1}
	\begin{center}
		\includegraphics[height=1.25cm]{media/alligator_red}
		\\
		\tikz[remember picture] \node (a) {
			\includegraphics[height=.75cm]{media/egg_yellow}
		};
		\begin{tikzpicture}[remember picture,overlay]
			\path (a.south) ++(0,-1) node[draw,anchor=east,ellipse callout,fill=white!50,opacity=1., callout absolute pointer={(a.south)},align=center]  {KRACK!\\[3mm] KRACK!};
		\end{tikzpicture}
	\end{center}
\end{frame}

\begin{frame}{Fressregeln: Beispiel 1}
	\begin{center}
		\includegraphics[height=1.25cm]{media/alligator_red}
		\\
		\includegraphics[height=.75cm]{media/egg_red}
	\end{center}
\end{frame}

\begin{frame}{Fressregeln: Beispiel 2}
	\begin{center}
		\begin{tabular}{rcccl}
			\multicolumn{2}{c}{
				\raisebox{-.5\height}{\includegraphics[height=1.25cm]{media/alligator_green}}
				\tikz[remember picture] \node (a) {\vphantom{X}};
				\begin{tikzpicture}[remember picture,overlay]
					\path<2> (a.east) ++(0,1) node[draw,anchor=west,ellipse callout,fill=white!50,opacity=1., callout absolute pointer={(a.mid)}]  {RAWR!};
				\end{tikzpicture}
			}
			& &
			\multicolumn{2}{c}{
				\raisebox{-.5\height}{\includegraphics[height=1.25cm]{media/alligator_green}}
				\tikz[remember picture] \node (b) {\vphantom{X}};
				\begin{tikzpicture}[remember picture,overlay]
					\path<2> (b.east) ++(0,1) node[draw,anchor=west,ellipse callout,fill=white!50,opacity=1., callout absolute pointer={(b.mid)}]  {HILFE!};
				\end{tikzpicture}
			}
			\\
			\raisebox{-.5\height}{\includegraphics[height=.75cm]{media/egg_green}}
			&
			\raisebox{-.5\height}{\includegraphics[height=.75cm]{media/egg_green}}
			& &
			\raisebox{-.5\height}{\includegraphics[height=.75cm]{media/egg_green}}
			&
			\raisebox{-.5\height}{\includegraphics[height=.75cm]{media/egg_green}}
		\end{tabular}
	\end{center}
\end{frame}

\begin{frame}{Fressregeln: Beispiel 2}
	\begin{center}
		\begin{picture}(180,80)
			\put(0, -20){
				\begin{turn}{7}
					\tikz[remember picture] \node (a) {\includegraphics[height=1.25cm]{media/alligator_green}};
				\end{turn}
			}
			\put(75, 7){
				\begin{turn}{15}
					\parbox{2cm}{
						\tikz[remember picture] \node (b) {\includegraphics[height=1.25cm]{media/alligator_green}};
						\\
						\hspace*{0.4cm}
						\includegraphics[height=.75cm]{media/egg_green}
						\hspace*{0.3cm}
						\includegraphics[height=.75cm]{media/egg_green}
					}
				\end{turn}
			}
		\end{picture}
		\begin{tikzpicture}[remember picture,overlay]
			\path (a.north) ++(1,1.5) node[draw,anchor=west,ellipse callout,fill=white!50,opacity=1., callout absolute pointer={(a.east)+(-0.2,1)}]  {YUMMY!};
		\end{tikzpicture}
		\begin{tikzpicture}[remember picture,overlay]
			\path (b.east) ++(0,1.5) node[draw,anchor=west,ellipse callout,fill=white!50,opacity=1., callout absolute pointer={(b.east)+(0,1)}]  {AHHHH!};
		\end{tikzpicture}
	\end{center}
	\vspace{.1cm}
	\hspace{2.8cm}
	\includegraphics[height=.75cm]{media/egg_green}
	\hspace{.6cm}
	\includegraphics[height=.75cm]{media/egg_green}
\end{frame}


\begin{frame}{Fressregeln: Beispiel 2}
	\begin{center}
		\begin{tabular}{cc}
			\multicolumn{2}{c}{
				\tikz[remember picture] \node (a) {\vphantom{X}};
				\begin{turn}{180}
					\raisebox{-.5\height}{\includegraphics[height=1.25cm]{media/alligator_green}}
				\end{turn}
				\begin{tikzpicture}[remember picture,overlay]
					\path (a.north) ++(1,1) node[draw,anchor=east,ellipse callout,fill=white!50,opacity=1., callout absolute pointer={(a.mid)}]  {ARGHHH};
				\end{tikzpicture}
			}

			\\
			\includegraphics[height=.75cm]{media/egg_green}
			&
			\includegraphics[height=.75cm]{media/egg_green}
		\end{tabular}
	\end{center}
\end{frame}

\begin{frame}{Fressregeln: Beispiel 2}
	\begin{center}
		\begin{tabular}{rcccl}
			\multicolumn{2}{c}{
				\raisebox{-.5\height}{\includegraphics[height=.75cm]{media/egg_green}}
				\tikz[remember picture] \node (a) {\vphantom{X}};
				\begin{tikzpicture}[remember picture,overlay]
					\path (a.east) ++(0,1) node[draw,anchor=west,ellipse callout,fill=white!50,opacity=1., callout absolute pointer={(a.mid)}]  {KRACK!};
				\end{tikzpicture}
			}
			& &
			\multicolumn{2}{c}{
				\raisebox{-.5\height}{\includegraphics[height=.75cm]{media/egg_green}}
				\tikz[remember picture] \node (b) {\vphantom{X}};
				\begin{tikzpicture}[remember picture,overlay]
					\path (b.east) ++(0,1) node[draw,anchor=west,ellipse callout,fill=white!50,opacity=1., callout absolute pointer={(b.mid)}]  {KRACK!};
				\end{tikzpicture}
			}

		\end{tabular}
	\end{center}
\end{frame}

\begin{frame}{Fressregeln: Beispiel 2}
	\begin{center}
		\begin{tabular}{rcccl}
			\multicolumn{2}{c}{
				\raisebox{-.5\height}{\includegraphics[height=1.25cm]{media/alligator_green}}
			}
			& &
			\multicolumn{2}{c}{
				\raisebox{-.5\height}{\includegraphics[height=1.25cm]{media/alligator_green}}
			}
			\\
			\raisebox{-.5\height}{\includegraphics[height=.75cm]{media/egg_green}}
			&
			\raisebox{-.5\height}{\includegraphics[height=.75cm]{media/egg_green}}
			& &
			\raisebox{-.5\height}{\includegraphics[height=.75cm]{media/egg_green}}
			&
			\raisebox{-.5\height}{\includegraphics[height=.75cm]{media/egg_green}}

		\end{tabular}
	\end{center}
\end{frame}

\begin{frame}{Spiel}
	\begin{center}
		\begin{tabular}{rccccl}
			\multicolumn{2}{c}{
				\includegraphics[height=1.25cm]{media/alligator_blank}
			}
			& &
			\includegraphics[height=1.25cm]{media/alligator_blank}
			& &
			\includegraphics[height=.75cm]{media/egg_pink}
			\\
			\includegraphics[height=.75cm]{media/egg_blank}
			&
			\includegraphics[height=.75cm]{media/egg_blank}
			& &
			\includegraphics[height=.75cm]{media/egg_blank}
		\end{tabular}
	\end{center}
\end{frame}

\section{Demo}

\begin{frame}
	\begin{center}
		\huge{DEMO}
	\end{center}
\end{frame}

\section{Der $\lambda$-Kalkül}

\begin{frame}[<+->]{Der $\lambda$-Kalkül}
	\begin{itemize}
		\item Programmierlogik
		\item Einführung bereits in den 1930ern
		\item Weite Verbreitung in der Lehre (KIT 5. Semester)
		\item Beispiel: $(\lambda x.\lambda y.x) (\lambda y.y)$
	\end{itemize}
\end{frame}

\begin{frame}[<+->]{Der $\lambda$-Kalkül}
	Visualisierung von Bret Victor:
	\begin{center}
		\begin{tabular}{rcl}
			\raisebox{-.5\height}{\includegraphics[height=1.cm]{media/egg_blank}}
			& $\Leftrightarrow$ &
			$x$
			\\[1cm]
			\raisebox{-.5\height}{\includegraphics[height=1.cm]{media/alligator_blank}}
			& $\Leftrightarrow$ &
			$\lambda x.$
		\end{tabular}
	\end{center}
\end{frame}

\begin{frame}{Die $\beta$-Reduktion}
	\begin{center}
		\begin{tabular}{lccl}
			&
			\raisebox{-.5\height}{
				\begin{tabular}{cc}
					\tikz[remember picture] \node(a) {
						\includegraphics[height=1.cm]{media/alligator_green}
					};
					&
					\tikz[remember picture] \node(b) {
						\includegraphics[height=.75cm]{media/egg_yellow}
					};
					\\
					\includegraphics[height=.75cm]{media/egg_green}
				\end{tabular}
				\begin{tikzpicture}[remember picture,overlay]
					\draw<2->[<-,thick,bend left=45] 
						($(a.east)+(-0.5,-0.1)$) to node[auto, swap] {}($(b.west)+(0.1,0)$);
				\end{tikzpicture}
			}
			& $\Leftrightarrow$ &
			$(\lambda x.x) y$
			
			\\[2cm]
			% SLIDE 3
			\action<3->{
				$\Rightarrow$
				&
				\visible<3->{
					\raisebox{-.5\height}{
						\parbox{3cm}{
							\begin{center}
								\begin{turn}{180}
									\includegraphics[height=1.cm]{media/alligator_green}
								\end{turn}
								\\
								\includegraphics[height=.75cm]{media/egg_yellow}
							\end{center}
						}
					}
				}
				& $\Leftrightarrow$ &
				$\cancel{\lambda x.}y$
			}
			
			\\[2cm]
			% SLIDE 4
			\action<4->{
				$\Rightarrow$
				&
				\visible<4->{\includegraphics[height=.75cm]{media/egg_yellow}}
				& $\Leftrightarrow$ &
				$y$
			}
		\end{tabular}
	\end{center}
\end{frame}

\section{Projektüberblick}

\begin{frame}

\end{frame}


\end{document}
